\subsection{The Role of Data Warehousing and DSS in Revenue Management}

\paragraph{From transactions to analysis.}
The intellectual premise of this project is stated by the textbook itself: ``\emph{increasingly, the success of a destination is determined by its ability to assemble, interpret, and utilize information in an effective manner}'' (p.~18)~\cite{goeldner2009}. A hotel's property management system (PMS) records booking transactions as they occur, but it is an operational system: it is organised to register each booking, not to answer analytical questions such as how occupancy and ADR evolved by market segment over two years. A data warehouse fills this role --- it \emph{remembers}: it integrates historical transactions and reorganises them into dimensional facts that can be rolled up, drilled down, sliced, and diced along the analytical dimensions~\cite{kimball2013}.
The integration problem is not hypothetical: recent reviews report that privacy-driven policies against sharing guest data have produced ``\emph{siloed storage for each property, leading to the unavailability of a 360-degree view of the customer}''~\cite{jawabreh2026} --- precisely the condition a warehouse exists to remedy.

\paragraph{Revenue management runs on decision support systems.}
In commercial practice, revenue management is executed by decision support systems, not by hand. State-of-the-art systems perform explicit price optimisation, overbooking, and allocation in sequence, fed by demand forecasting~\cite{pimentel2019}, and their adoption is associated with revenue gains of approximately 6\%~\cite{boyd2003} (American Airlines' pioneering system reported 4--7\%~\cite{smith1992}). The three decisions are quantitative at their core: forecasting demand by segment, setting protection levels between fare classes (the two-class closed form dates to Littlewood~\cite{littlewood2005}), valuing capacity through bid prices, and balancing oversell against undersell costs. Running these algorithms on real data --- rather than on intuition or spreadsheets, whose ``success is largely dependent on the caliber of the person''~\cite{anderson2012} --- is precisely what a decision support system provides. Turban et al.~characterise a DSS as combining a data management subsystem, a model management subsystem, a user interface, and a knowledge-based component~\cite{turban2011}; our architecture (Section~3.1) follows this decomposition.

\paragraph{The data foundation.}
The project is grounded in real data: the hotel booking demand dataset of Antonio et al.~\cite{antonio2019}, extracted directly from the PMS databases of two Portuguese hotels (a resort in the Algarve and a city hotel in Lisbon) and described fully in Section~2. Beyond its industry specificity, the dataset was constructed with an anti-leakage principle that matches the needs of decision models: variable values were recorded as of the day prior to each arrival, so each record is an on-the-books snapshot rather than a post-hoc summary --- the same temporal discipline our overbooking model adopts (Section~3.3).
The dataset has itself already supported a published decision support system for cancellation prediction~\cite{antonio2019dss}, and a survey of public alternatives found no hotel dataset richer in revenue-management variables: every public ``hotel booking'' dataset traces back to it, and its direct successor --- four hotels, 2014--2023, 670{,}343 bookings --- is available only under NDA~\cite{silvestre2025}.

\paragraph{The course-adopted topic: the MBMS.}
The course-adopted topic for this group is the \emph{Model Base Management Subsystem} (MBMS) --- the component that manages, manipulates, and runs the models embedded in the DSS. Turban et al.~describe it as ``a module that theoretically has capabilities similar to those of a DBMS to manage, manipulate, and run the models embedded in the DSS''~\cite{turban2011}, and no standard, off-the-shelf MBMS exists: every DSS team designs its own~\cite{turban2011}. Our design --- presented in depth in Section~3.4 --- implements the model base as database tables: a \texttt{model\_catalog} of model metadata, a uniform execution registry, and a scenario run log.
